% A sample beamer presentation showing 
% what the UNCMath theme looks like.
%
% Key features of the template:
%
% - Little excess colors
% - Using cleaner itemize/enumerate (no "spheres" or "balls")
% - Simple color scheme; shades of "Carolina blue", grayscale
% - small UNC logo in upper right of every slide.
% - removal of some of the ugly default beamer things:
% -- navigation menus
% -- date in footer
% -- page numbering in footer
%
% This is all with the purpose of improving the 
% "usable" space with beamer, and removing 
% the doodads that most people never 
% use or look at when they present or 
% watch someone with a beamer presentation.
%
% First version created by Manuchehr Aminian, April 2017.
% Feel free to modify and share, make suggestions.
% 
% There's no need to credit me, but it would 
% make me happy to know that people used it. 
% Feel free to buy me food or drinks, though.
%

\documentclass{beamer}

\usepackage{graphicx}
%\usepackage{epstopdf}

% Import the UNCMath theme.
\usetheme{UNCMath}

\logo{
\includegraphics[height=3em]{oldwell_cmyk}
}

% A macro for making the title frame.
% Removes the bottom bar and logo temporarily.
% If you don't want these in other frames, 
% you could try mimicking this.
\newcommand{\titleframe}{
{
\setbeamertemplate{headline}{}
\setbeamertemplate{footline}{}
\begin{frame}
\titlepage
\end{frame}
}
}


%%%%%%%%%%%%%%%%%%%%%%%%%%%%%%%%%%%%%%%%%%
%
% Information for the title frame.
%
\title{A Very Important Mathematical Problem}
\subtitle{or, A Big Uncertain Data Quantification}
\author[Manuchehr Aminian]{Manuchehr Aminian}
\institute{University of North Carolina, Chapel Hill}
\date{April 6th, 2017}

%%%%%%%%%%%%%%%%%%%%%%%%%%%%%%%%%%%%%%%%%%
%
% Start of document.
%
\begin{document}

% This is the titleframe, do it this way 
% so that you don't have the logo on the title page.
\titleframe


\begin{frame}
\frametitle{Titular Statement}
It is a trivial matter to demonstrate that 
%
\begin{equation}
e^{i\pi}+1=0.
\end{equation}
%
Here is a ``sketch" of the proof:
\begin{enumerate}
\item First, assume that $e^{i \theta} = \cos(\theta) + i \sin(\theta)$.
\begin{itemize}
\item (any third grader knows this)
\end{itemize}
\item Then, substitute $\theta = \pi$, and use your knowledge of trigonometric functions.
\begin{enumerate}
\item (this was done in second grade)
\item (I'm really just testing how itemize and enumerate interact)
\end{enumerate}
\end{enumerate}
%
Between lists of things, some more text is added.
%
\begin{itemize}
\item Test this out too
\begin{enumerate}
\item testing yet another interaction
\end{enumerate}
\item No idea what will happen
\begin{itemize}
\item some further unordered points go here
\end{itemize}
\end{itemize}

\end{frame}
%
\begin{frame}
\frametitle{Here is a very long frame title which 
will extend on to a second line. }
Notice that the title does not run through 
the logo. This took me a lot of suffering to sort out 
(along with the logo positioning/behavior itself). 
The documentation isn't that great out there.
\vspace{1em}

Another small annoyance is that ``vspace"s are needed 
to have spaces between pieces of text.
\vspace{1em}

With all this in mind, we state a theorem.
\begin{theorem}[A trivial statement]
Suppose that $e^{i\theta} = \cos(\theta) + i\sin(\theta)$, 
where $i=\sqrt{-1}$ and $\theta \in \mathbb{R}$. \\[1em]
Then $e^{i\pi}+1=0$.
\end{theorem}
\end{frame}
%
\begin{frame}
If you don't include a frame title, there won't be a horizontal 
rule included, but the logo will still show. 
If you don't want the logo in a frame like this 
either, you need to mimick the behavior done with the 
title page macro defined in the source file.
\end{frame}
%
\begin{frame}
\frametitle{Here's the UNC logo}

\begin{itemize}
\item I've included the Old Well logo here, 
because we all love our university very much. 
\item Keep in mind 
that you \textbf{need} to have the file \texttt{oldwell\_cmyk.eps} 
in the same directory as your \texttt{.tex} file for the 
presentation to compile successfully. In principle, 
you could use a different image.
\item The Old Well logo is from 
UNC's ``Downloadable Resources:"
{\color{blue!80} \url{http://identity.unc.edu/downloadable-resources/}}
\end{itemize}

\begin{figure}
\centering
\includegraphics[width=0.6\linewidth]{oldwell_cmyk}
\end{figure}
\end{frame}
%

\end{document}
